% Project References Assignment -- TIF360/FYM360
% Compile on Overleaf with: project_references_sarcasm.tex + references_sarcasm.bib

\documentclass[%
 reprint,
 amsmath,amssymb,
 aps,
]{revtex4-2}

\usepackage{graphicx}
\usepackage{dcolumn}
\usepackage{bm}
\usepackage{xcolor}
\usepackage{hyperref}
\usepackage{enumitem}
\usepackage{microtype}

\begin{document}

\title{[C] Context-Aware and Retrieval-Augmented Sarcasm Detection in Social Media}

\author{Mina Tahmasebi Berjouei}

\date{\today}

\begin{abstract}
This document submits the project references for the project \emph{Context-Aware and Retrieval-Augmented Sarcasm Detection in Social Media}. The selected references support the main components of the project: the Reddit-based sarcasm dataset, context-aware sarcasm modelling, transformer-based supervised classifiers, large-language-model baselines, retrieval-augmented classification, and interpretability. For each reference, I briefly explain why it is relevant to the project and how it supports the planned methodology.

\begin{description}
\item[Project Topic]
{Sarcasm Detection in Social Media.}
\item[Teaching Assistant]
{Antonio Ciarlo.}
\end{description}
\end{abstract}

\maketitle

\section{Project References and Relevance}

The project investigates whether a supervised transformer classifier that uses both a target Reddit comment and its conversational context can detect sarcasm more effectively than context-free baselines, and whether retrieval-augmented examples or LLM-based reasoning can improve interpretability. The following references are selected because they cover the required dataset foundation, core NLP models, context-aware sarcasm detection methods, retrieval-augmented methods, and explanation-oriented approaches.

\begin{enumerate}[leftmargin=*, itemsep=0.65em]

\item \textbf{Khodak, Saunshi, and Vodrahalli -- SARC dataset \cite{khodak2018large}.} This is the most important dataset reference for the project. It introduces the Self-Annotated Reddit Corpus (SARC), a large Reddit sarcasm corpus with author-provided sarcasm labels, parent-comment context, subreddit information, and author metadata. It directly justifies the dataset choice and supports experiments comparing target-only and context-aware models.

\item \textbf{Oprea and Magdy -- intended sarcasm \cite{oprea2020isarcasm}.} This paper is relevant because it distinguishes intended sarcasm from sarcasm perceived by external annotators. It supports the decision to use self-annotated data, since sarcasm is subjective and third-party annotators may miss the author's intended meaning.

\item \textbf{Wallace et al. -- humans require context \cite{wallace2014humans}.} This reference provides an early and important motivation for context-aware sarcasm detection. It shows that even human readers need contextual information to infer ironic intent, which supports the project's central hypothesis that models should use parent-comment context.

\item \textbf{Riloff et al. -- sentiment-situation contrast \cite{riloff2013sarcasm}.} This paper is relevant because it frames sarcasm as a contrast between positive surface sentiment and a negative situation. That idea directly supports the project's focus on incongruity between the reply and the conversational context.

\item \textbf{Joshi, Sharma, and Bhattacharyya -- context incongruity \cite{joshi2015harnessing}.} This work is relevant because it explicitly uses context incongruity as a signal for sarcasm. It provides theoretical support for designing features or attention mechanisms that detect mismatch between literal wording and surrounding context.

\item \textbf{Hazarika et al. -- CASCADE \cite{hazarika2018cascade}.} CASCADE is a key context-aware sarcasm detection framework. It incorporates conversational and user-level context, making it useful for motivating the broader view that sarcasm depends not only on the target utterance but also on discourse and speaker behaviour.

\item \textbf{Dong et al. -- transformer-based context-aware sarcasm detection \cite{dong2020transformer}.} This is one of the closest methodological references to the project. It evaluates transformer-based sarcasm detection in conversation threads and directly supports the planned input format where the model receives both the previous context and the target reply.

\item \textbf{Baruah, Das, Barbhuiya, and Dey -- BERT with conversational context \cite{baruah2020context}.} This paper is relevant because it studies BERT-based sarcasm detection with contextual information. It also motivates careful ablation, since context can help but may introduce noise if it is not integrated properly.

\item \textbf{Devlin et al. -- BERT \cite{devlin2019bert}.} BERT is a foundational transformer model for contextual text representation. It is relevant as a strong neural baseline and as the conceptual basis for fine-tuning transformer classifiers on sarcasm detection.

\item \textbf{Liu et al. -- RoBERTa \cite{liu2019roberta}.} RoBERTa is relevant because it improves BERT pretraining and is commonly used as a strong supervised NLP classifier. It is a suitable candidate for the project's context-free and context-aware transformer baselines.

\item \textbf{He et al. -- DeBERTa \cite{he2021deberta}.} DeBERTa is relevant as a stronger transformer candidate because its disentangled attention separately models content and position. This is useful for paired parent-comment/reply inputs, where the model must learn cross-sentence relations and pragmatic incongruity.

\item \textbf{Potamias, Siolas, and Stafylopatis -- RCNN-RoBERTa \cite{potamias2020roberta}.} This work is relevant because it applies RoBERTa-based representations to sarcasm detection and provides a strong deep-learning benchmark. It supports the use of RoBERTa as a serious baseline rather than only classical machine-learning methods.

\item \textbf{Lewis et al. -- Retrieval-Augmented Generation \cite{lewis2020retrieval}.} This is the foundational RAG reference. It is relevant because the project explores retrieval-augmented examples as an additional evidence source for classification and explanation.

\item \textbf{Wu et al. -- RAG survey \cite{wu2024rag}.} This survey is useful for understanding the technical design of a RAG pipeline, including retrievers, external knowledge stores, and retrieval-generation integration. It supports the implementation design of the retrieval component.

\item \textbf{Yu et al. -- retrieval-augmented few-shot text classification \cite{yu2023retrieval}.} This paper is relevant because it adapts retrieval augmentation to classification rather than only question answering. It supports the idea of retrieving similar labelled examples and using them as evidence for sarcasm prediction.

\item \textbf{Zhang et al. -- SarcasmBench \cite{zhang2024sarcasmbench}.} This reference is relevant for evaluating LLMs on sarcasm understanding. It supports the project's decision to include LLM zero-shot or few-shot classification as a comparison method, while still keeping supervised transformers as the main model.

\item \textbf{Yao et al. -- SarcasmCue \cite{yao2025sarcasmcue}.} This work is relevant because it studies whether sarcasm detection in LLMs should be treated as step-by-step reasoning. It helps design fair LLM prompting baselines and cautions against assuming that chain-of-thought prompting always improves sarcasm detection.

\item \textbf{Lee et al. -- pragmatic metacognitive prompting \cite{lee2025pragmatic}.} This paper is relevant because it uses linguistic-pragmatic theories to improve LLM sarcasm detection. It can guide the design of explanation prompts that ask an LLM to consider implicature, pretense, and speaker intent.

\item \textbf{Xiong et al. -- Sarc7 \cite{xiong2025sarc7}.} This reference is relevant because it evaluates sarcasm across multiple sarcasm types and emphasizes emotion-informed techniques. It supports the project's planned explanation analysis, especially when interpreting emotional contrast between context and reply.

\item \textbf{Yang et al. -- Sarcasm-R1 \cite{yang2025sarcasmr1}.} This paper is relevant because it decomposes sarcasm reasoning into language, context, and emotion. That structure is useful for designing qualitative explanations and for analysing why a model predicts sarcasm.

\item \textbf{Qiu et al. -- commonsense reasoning for emotional incongruity \cite{qiu2025emotional}.} This work is relevant because sarcasm often requires commonsense knowledge and emotional incongruity detection. It supports a possible extension where retrieval provides not only similar examples but also commonsense evidence.

\item \textbf{Wen and Rezapour -- prototype-based interpretable sarcasm model \cite{wen2026prototype}.} This reference is relevant because it focuses on interpretable contextual sarcasm detection using transformer and prototype-based representations. It supports the project's goal of making predictions more explainable through similar examples or prototype-like evidence.

\end{enumerate}

\section{Summary of Selected Literature}

The selected literature supports a coherent project design. First, SARC and iSarcasm justify the use of self-annotated sarcasm data rather than relying only on third-party annotation \cite{khodak2018large,oprea2020isarcasm}. Second, earlier sarcasm studies motivate the importance of context and incongruity \cite{wallace2014humans,riloff2013sarcasm,joshi2015harnessing}. Third, context-aware neural models and transformer methods justify the main supervised modelling approach \cite{hazarika2018cascade,dong2020transformer,baruah2020context,devlin2019bert,liu2019roberta,he2021deberta}. Finally, RAG, LLM evaluation, and interpretable/prototype-based methods support the additional comparison and explanation components \cite{lewis2020retrieval,wu2024rag,yu2023retrieval,zhang2024sarcasmbench,yao2025sarcasmcue,lee2025pragmatic,xiong2025sarc7,yang2025sarcasmr1,qiu2025emotional,wen2026prototype}. Therefore, these references are directly aligned with the project question and provide enough coverage for dataset choice, baseline design, main model selection, retrieval augmentation, and interpretability.

\bibliographystyle{apsrev4-2}
\bibliography{references_sarcasm}

\end{document}
